\chapter*{INTRODUCTION GÉNÉRALE}
\markboth{\MakeUppercase{INTRODUCTION GÉNÉRALE}}{}
%\addstarredchapter{INTRODUCTION GÉNÉRALES}
\addcontentsline{toc}{chapter}{INTRODUCTION GÉNÉRALE}
\adjustmtc
\thispagestyle{MyStyle}


L'avènement de l'intelligence artificielle a ouvert de nouvelles perspectives dans le domaine de l'accessibilité numérique. Parmi les applications prometteuses de cette technologie figure la conversion d'images en braille en texte lisible, un outil essentiel pour faciliter la communication et l'accès à l'information pour les personnes malvoyantes.

Ce rapport explore le processus de développement d'un modèle d'IA dédié à cette tâche, mettant en lumière les défis techniques rencontrés, les solutions implémentées et les bénéfices potentiels pour les utilisateurs finaux. En tant que stagiaire, mon rôle a été de contribuer à la conception et à la réalisation de ce modèle, en utilisant des technologies modernes telles que Python, TensorFlow et des datasets.

Nous commencerons par examiner le contexte et la nécessité croissante de tels outils de conversion dans le cadre de l'inclusion numérique. Ensuite, nous plongerons dans les détails techniques du développement, en mettant en évidence les choix architecturaux, les fonctionnalités clés implémentées et les considérations de précision et de performance du modèle.

Enfin, nous évaluerons l'impact potentiel de cette application sur les utilisateurs malvoyants, en analysant les résultats obtenus et les perspectives d'amélioration future. Ce rapport vise à offrir une compréhension approfondie du processus de développement d'un modèle d'IA pour la conversion d'images braille en texte, ainsi qu'une vision claire de son rôle dans l'amélioration de l'accessibilité numérique et de l'expérience utilisateur pour les personnes malvoyantes.\par

